\section{Surface estimates and consequences}\label{sec:surfaces}

\subsection{Passage to the surface}

We now deduce Theorem~\ref{thm:main} from
Theorem~\ref{thm:graph}. Cover \(\Sigma\) by finitely many sufficiently
small charts. By Lemma~\ref{lem:continuation}, after fixed affine
normalizations each chart lies in a graph with a phase admissible
for Theorem~\ref{thm:packet}. The construction and its constants are
fixed separately for each chart and do not depend on \(p\).
A measurable partition subordinate to the covering avoids any
assumption that the data extend smoothly across chart boundaries.

Write one chart in ambient frequency coordinates as
\[
 \xi=v_0+L(\zeta,\phi(\zeta)),\qquad \zeta\in A\subset\R^2,
\]
where \(L\) is invertible and \(A\) is bounded, possibly cut off by
the boundary of the surface. Its surface measure is
\(J_\Sigma(\zeta)\dd\zeta\), with \(J_\Sigma\) smooth and strictly
positive on a neighborhood of the chart. For data supported on
this piece, set
\[
 g(\zeta)=
 \begin{cases}
  f\bigl(v_0+L(\zeta,\phi(\zeta))\bigr)J_\Sigma(\zeta),
       &\zeta\in A,\\
  0,&\zeta\notin A.
 \end{cases}
\]
Then
\[
 E_\Sigma f(x)
 =e(x\cdot v_0)\int_{\R^2}g(\zeta)
       e\bigl((L^{\mathsf T}x)'\cdot\zeta+
                (L^{\mathsf T}x)_3\phi(\zeta)\bigr)\dd\zeta.
\]
The change of physical variables \(y=L^{\mathsf T}x\) contributes
\(|\det L|^{-1/p}\) to the \(L^p\) norm. Moreover,
\begin{align*}
 \|g\|_3^3
 &\le(\sup_A J_\Sigma)^2
       \int_{\Sigma_{\mathrm{chart}}}|f|^3\dd\sigma,\\
 \|f\|_{L^3(\Sigma_{\mathrm{chart}})}
 &\le
 \sigma(\Sigma_{\mathrm{chart}})^{1/3-1/p}
 \|f\|_{L^p(\Sigma_{\mathrm{chart}})},\qquad p>3.
\end{align*}
After the finite window partition of
Equation~\eqref{eq:window-partition}, Theorem~\ref{thm:graph} applies
to every resulting term. The determinant, area, window, and chart
factors are bounded uniformly for \(3<p\le4\). Summing the finitely
many terms by the triangle inequality gives
Equation~\eqref{eq:main} in this range, with the quantitative bound
Equation~\eqref{eq:quantitative-intro}. Here \(P_\Sigma\) is the
maximum of the finitely many packet constants, enlarged to be at
least one.

For \(p>4\), the local \(L^4\) and \(L^\infty\) bounds give the result
directly, as in Theorem~\ref{thm:graph}, and the same finite chart
argument applies. No smoothness of \(f\), nor of its extension by
zero on the parameter plane, has been used. Boundary charts are
therefore included.

Proposition~\ref{prop:lower} supplies
Equation~\eqref{eq:lower-intro} for every nonempty surface by choosing
a nonzero smooth datum in an interior chart. This completes the
proof of Theorem~\ref{thm:main} and the stated dependence on \(p\).

\subsection{Free Schr\"odinger local smoothing}

The extension estimate for a compact paraboloid controls all frequencies
of the Schr\"odinger propagator through the transfer theorem of
Lee, Rogers, and Seeger. The strict Sobolev inequality in
Corollary~\ref{cor:schrodinger-smoothing} then supplies the summability
across frequency scales.

\begin{proof}[Proof of Corollary~\ref{cor:schrodinger-smoothing}]
On the compact paraboloid
\(\Sigma=\{(y,|y|^2):|y|\le1\}\), write
\(\dd\sigma=w(y)\dd y\), where \(w(y)=\sqrt{1+4|y|^2}\).
Apply Theorem~\ref{thm:main} to the datum \(g(y)/w(y)\) on \(\Sigma\).
A fixed change of output variables then gives
\[
 \left\|\int_{|y|\le1}g(y)e^{i(\tau|y|^2-\xi\cdot y)}\dd y\right\|_{L^p(\R^3_{\xi,\tau})}
 \le C_p\|g\|_{L^p(\{|y|\le1\})}.
\]
This is the compact-frequency extension estimate in
Lee--Rogers--Seeger \cite[Equation~(1.3)]{LeeRogersSeeger2013} with
both exponents equal to \(p\). Their Theorem~1.1, with spatial dimension
two and Besov summation exponent two, yields
\[
 \|e^{it\Delta}f\|_{L^p(\R^2\times[0,1])}
 \le C_p\|f\|_{B^p_{\alpha,2}(\R^2)},
 \qquad \alpha=2(1-2/p)-2/p=2-6/p.
\]
For an inhomogeneous Littlewood--Paley partition \((P_k)_{k\ge0}\),
the uniform annular multiplier bound
\(\|P_kf\|_p\le C_s2^{-ks}\|(I-\Delta)^{s/2}f\|_p\) gives
\[
 \|f\|_{B^p_{\alpha,2}}
 \le C_s\left(\sum_{k\ge0}2^{-2k(s-\alpha)}\right)^{1/2}
       \|f\|_{W^{s,p}}.
\]
The series converges for \(s>\alpha\). This proves
Equation~\eqref{eq:schrodinger-smoothing} on Schwartz functions, and
density gives the stated bounded extension.
\end{proof}

\subsection{Translation-invariant oscillatory integrals}

A second application concerns phases whose curvature changes only by
a scalar factor in time. Such a phase can be straightened by changing
the output variables alone, reducing the estimate to one fixed
elliptic graph. This is the translation-invariant setting of the output
straightening theorem of Gao--Liu--Xi
\cite[Theorem~1.12 and Section~4.1]{GaoLiuXi2025}.
We give the short argument for the two-dimensional input considered
here, then allow a general smooth amplitude by Fourier expansion.

\begin{corollary}[Elliptic translation-invariant Bourgain phases]
\label{cor:bourgain-oscillatory}
Let
\[
 \phi(x,t;y)=x\cdot y+\psi(t;y),\qquad
 x,y\in\R^2,\quad t\in\R,\quad \psi(0;y)=0,
\]
be a fixed real smooth phase near the origin. Suppose throughout
this neighborhood that
\[
 M(t,y)=D_y^2\partial_t\psi(t;y)\quad\hbox{is definite},
 \qquad \partial_tM(t,y)=c(t,y)M(t,y)
\]
for a smooth real scalar function \(c\). On a sufficiently small
fixed product chart, let \(a\) be any fixed smooth amplitude with
compact support in its interior. For every \(3<p<\infty\),
\begin{equation}\label{eq:bourgain-oscillatory}
 \left\|\int e^{iN\phi(x,t;y)}a(x,t;y)f(y)\dd y
       \right\|_{L^p(\R^3_{x,t})}
 \le C_{\phi,a,p}N^{-3/p}\|f\|_{L^p(\R^2)},\qquad N\ge1,
\end{equation}
where the operator is zero outside the output support of \(a\).
The constant may depend on the fixed chart, phase, amplitude and
\(p\), but not on \(N\) or \(f\).
\end{corollary}

\begin{proof}
Shrink the product neighborhood so that its input domain is a ball.
Since \(M\) is a Hessian in \(y\), commuting its derivatives in
\(\partial_tM=cM\) gives
\[
 (\partial_{y_k}c)M_{ij}=(\partial_{y_i}c)M_{kj}.
\]
Multiplication by \((M^{-1})_{j\ell}\), summed over \(j\), yields
\((\partial_{y_k}c)\delta_{i\ell}
 =(\partial_{y_i}c)\delta_{k\ell}\).
For each \(k\), choose \(\ell=i\ne k\); hence
\(\nabla_yc=0\). The input ball is connected, so \(c=c(t)\).
Set
\[
 \rho(t)=\exp\!\left(\int_0^t c(v)\,dv\right),\qquad
 h(y)=\partial_t\psi(0;y),\qquad
 S(t)=\int_0^t\rho(v)\,dv.
\]
The equation for \(M\) now implies
\(D_y^2(\partial_t\psi-\rho h)=0\). Thus
\(\partial_t\psi=\rho(t)h(y)+b(t)\cdot y+d(t)\)
for smooth functions \(b,d\). Integrating and using \(\psi(0;y)=0\)
gives
\[
 \psi(t;y)=S(t)h(y)+B(t)\cdot y+D(t),\qquad
 B(t)=\int_0^t b(v)\,dv,\quad D(t)=\int_0^t d(v)\,dv.
\]
Since \(S'=\rho>0\), it has a smooth local inverse \(t=t(s)\).
The output change
\(\kappa(u,s)=(u-B(t(s)),t(s))\) therefore satisfies
\[
 \phi(\kappa(u,s);y)=u\cdot y+s h(y)+r(s),\qquad
 r(s)=D(t(s)).
\]
Here \(D_y^2h=M(0,y)\) is definite. No input coordinate change
is needed.

Put \(z=(u,s)\) and \(A(z,y)=a(\kappa(z);y)\). Extend this fixed
amplitude smoothly by zero, and periodically in \(z\) across a fixed
cube \(Q\) of side \(L\) whose interior contains its output support.
On \(Q\),
\[
 A(z,y)=\sum_{k\in\Z^3}e^{2\pi i k\cdot z/L}A_k(y),
 \qquad \sum_k\|A_k\|_\infty<\infty.
\]
Indeed, integration by parts in \(z\) gives
\(\|A_k\|_\infty\le C_m(1+|k|)^{-m}\) for any fixed \(m>3\).
All coefficients have a common compact frequency support, and the
constants are independent of \(N\).
Choose a compact graph patch \(\Sigma\), parametrized by
\(\omega(y)=(y,h(y))\), with smooth boundary and parameter domain
containing this support, and write
\(\dd\sigma=w(y)\dd y\), where \(w(y)=\sqrt{1+|\nabla h(y)|^2}\).
For
\[
 g_k(\omega(y))=A_k(y)f(y)/w(y)
\]
we have
\[
 \|g_k\|_{L^p(\Sigma)}^p
 =\int |A_kf|^p w^{1-p}\dd y
 \le\|A_k\|_\infty^p\|f\|_p^p.
\]
The factor \(e^{iNr(s)}\) is unimodular, and each remaining Fourier
term is \(e^{2\pi i k\cdot z/L}E_\Sigma g_k(Nz/(2\pi))\).
Minkowski's inequality and Theorem~\ref{thm:main}, applied globally
to each term, bound their sum on \(Q\) by
\[
 (2\pi/N)^{3/p}C_{\Sigma,p}
       \sum_k\|A_k\|_\infty\|f\|_p.
\]
Changing back by \(\kappa\) contributes only the bounded factor
\(\sup|\det D\kappa|^{1/p}\) on the fixed output support.
Absolute coefficient summability and the local \(L^1\) integrability
of \(f\) justify the series and its interchange with the integral.
This proves Equation~\eqref{eq:bourgain-oscillatory}, with no power
loss in \(N\).
\end{proof}
